\section{Статические игры с неполной информацией}

\begin{example}[Мотивирующий пример: Pepsico vs. Coca-Cola]
Две компании, Pepsico и Coca-Cola, планируют развиваться в новом сегменте рынка (новый тип энергетических напитков). Coca-Cola уже присутствует на этом рынке и может увеличить долю ($U$) или оставить всё как есть ($D$). Pepsico может войти на рынок ($In$) или не входить ($Out$). 

Фирмы принимают решение одновременно. Проблема заключается в том, что Pepsico точно не знает, насколько затратно для Coca-Cola расшириться. Стоимость увеличения доли может быть либо высокой ($High$) с вероятностью $\theta$, либо низкой ($Low$) с вероятностью $1 - \theta$.

\begin{table}
\centering
\begin{NiceTabular}{cc|cc}
\multicolumn{4}{c}{\textbf{Высокие издержки (High costs)}} \\
\toprule
& & \multicolumn{2}{c}{\textbf{Pepsico}} \\
& & \textit{In} & \textit{Out} \\
\midrule
\multirow{2}{*}{\textbf{Coca-Cola}} 
& $U$ & $0, -1$ & $2, 0$ \\
& $D$ & $2, 1$  & $3, 0$ \\
\bottomrule
\end{NiceTabular}
\quad
\begin{NiceTabular}{cc|cc}
\multicolumn{4}{c}{\textbf{Низкие издержки (Low costs)}} \\
\toprule
& & \multicolumn{2}{c}{\textbf{Pepsico}} \\
& & \textit{In} & \textit{Out} \\
\midrule
\multirow{2}{*}{\textbf{Coca-Cola}} 
& $U$ & $3, -1$ & $5, 0$ \\
& $D$ & $2, 1$  & $3, 0$ \\
\bottomrule
\end{NiceTabular}
\end{table}

Игрок 1 (Coca-Cola) знает тип своих издержек и может выбирать стратегию в зависимости от их величины (type-specifically). Анализ матриц позволяет определить строго доминирующие стратегии для каждого типа первого игрока:
\begin{align*}
    a_1^*(t = H) &= D, \\
    a_1^*(t = L) &= U.
\end{align*}

Игрок 2 (Pepsico) знает только распределение вероятностей $\theta$, но может рационально предвидеть поведение оппонента $a_1^*(t)$. При этом второй игрок должен выбрать единственную стратегию. Ожидаемая полезность Pepsico для каждой из стратегий:
\begin{align*}
    \mathrm{E}u_2(a_1^*(H), a_1^*(L), In) &= 1\cdot\theta + (-1)(1 - \theta) = -1 + 2\theta, \\
    \mathrm{E}u_2(a_1^*(H), a_1^*(L), Out) &= 0\cdot\theta + 0(1 - \theta) = 0.
\end{align*}

Оптимальная стратегия второго игрока зависит от значения параметра $\theta$:
\[
a_2^* = 
\begin{cases}
    In, & \text{если } \theta > 1/2 \\
    \alpha In + (1 - \alpha) Out, \: \alpha \in [0, 1], & \text{если } \theta = 1/2 \\
    Out, & \text{если } \theta < 1/2
\end{cases}
\]
\end{example}

\subsection{Формализация}

Стратегическая игра с неполной информацией называется \textit{байесовой игрой} (игрой с неполной, асимметричной информацией). Ключевой элемент описания такой игры — множество \textit{состояний мира}. Каждое состояние представляет собой полное описание одного набора характеристик игроков, включая особенности доступной им информации и функции полезности.

В начале игры реализуется конкретное состояние мира $\omega$, которое игроки напрямую не наблюдают. Вместо этого каждый игрок получает \textit{сигнал}, раскрывающий определенную информацию об истинном состоянии мира. 

Обозначим сигнал, который игрок $i$ получает в состоянии $\omega$, через $\tau_i(\omega)$. Функция $\tau_i$ является детерминированной функцией от состояния игры. Состояние игры, порождающее сигнал $t_i$, называется \textit{согласованным} с $t_i$. Размер множеств состояний, согласованных с сигналами каждого игрока, отражает качество доступной ему информации.

Полученный игроком $i$ сигнал $t_i$ формально называется \textit{типом} $t_i$ игрока $i$. Игрок каждого типа имеет определенные представления о вероятности реализации того или иного состояния, согласованного с его сигналом. Если сигнал согласован с единственным состоянием мира, игрок считает его верным с вероятностью $1$. Поскольку игрокам важны не только действия оппонентов, но и состояние мира, их предпочтения определяются над вероятностным распределением на парах $(a, \omega)$.

\begin{definition}[Байесова игра]
Байесова игра состоит из следующих элементов:
\begin{itemize}
    \item множество игроков $I = \{1, \dots, n\}$;
    \item для каждого игрока $i$, множество типов $T_i$ (или, эквивалентно, множество сигналов и представление о том, какие состояния мира согласованы с сигналом — распределение вероятностей на множестве состояний мира);
    \item для каждого игрока $i$, множество действий $A_i$;
    \item функция полезности, заданная на парах $(a, \omega)$.
\end{itemize}
\end{definition}

Пусть $T_i$ — множество типов игрока $i$. Множество профилей типов всех игроков определяется как $T = \times_{i=1}^N T_i$. Конкретный профиль из $T$ можно интерпретировать как состояние мира. Пусть $P$ — распределение вероятностей на $T$. 

Если игрок $i$ имеет тип $t_i \in T_i$ (и знает его), его представления о типах остальных игроков формируются посредством пересчета вероятностей по формуле Байеса:
\[
P(t_{-i} | t_i) = \frac{P(t_i, t_{-i})}{P(t_i)} = \frac{P(t_i, t_{-i})}{\sum_{t_{-i} \in T_{-i}} P(t_i, t_{-i})}
\]

Стратегия игрока $i$ является функцией, зависящей от его типа: $s_i: T_i \to A_i$. Ожидаемая полезность игрока $i$ типа $t_i$ вычисляется как:
\[
\tilde{u}_i(s, t_i) = \sum_{t_{-i} \in T_{-i}} P(t_{-i} | t_i) u_i(s_i(t_i), s_{-i}(t_{-i}), t_i).
\]

\begin{definition}[Равновесие Байеса-Нэша]
Пусть $\langle I, A, T, P, u \rangle$ — байесова игра. Равновесие Байеса-Нэша — это профиль стратегий $s^*$, такой что для каждого игрока $i$, для всех возможных альтернативных действий $a'_i \in A_i$, и для каждого типа $t_i \in T_i$ игрока $i$, выполняется условие:
\[
\sum_{t_{-i} \in T_{-i}} P(t_{-i} | t_i) u_i(s_i^*(t_i), s_{-i}^*(t_{-i}), t_i) \ge \sum_{t_{-i} \in T_{-i}} P(t_{-i} | t_i) u_i(a'_i, s_{-i}^*(t_{-i}), t_i).
\]
\end{definition}

\begin{remark}
Байесова игра эквивалентна игре в развернутой форме (с информационными множествами), которая начинается с хода Природы (игрока 0), случайным образом распределяющей типы между игроками. Отношение доминирования стратегий в таких играх определяется отдельно для каждого типа.
\end{remark}

\begin{theorem}
В любой байесовой игре с конечными множествами $I$, $A$ и $T$ существует равновесие Байеса-Нэша.
\end{theorem}

\subsection{Примеры статических игр с неполной информацией}

\begin{example}[Дуополия Курно с неизвестными издержками]
Две фирмы (1 и 2) производят однородный продукт и конкурируют по объему выпуска. Фирмы одновременно выбирают объемы производства $q_1$ и $q_2$. Единая рыночная цена формируется как $p = 1 - q_1 - q_2$. Функция полезности (прибыль) фирмы $i$:
\[
u_i = q_i(1 - q_1 - q_2) - c_i q_i
\]
Фирмы различаются по издержкам $c_1$ и $c_2$:
\begin{itemize}
    \item значение $c_1$ известно обеим фирмам;
    \item значение $c_2$ является частным знанием фирмы 2. Известно, что $c_2 = c_H$ с вероятностью $\theta \in (0, 1)$, и $c_2 = c_L$ (где $c_L < c_H$) с вероятностью $1 - \theta$. Это описание является общим знанием.
\end{itemize}

Множество типов фирмы 2: $T = \{H, L\}$. Равновесие Байеса-Нэша представляет собой тройку величин $(q_1^*, q_2^*(H), q_2^*(L))$, где:
\begin{itemize}
    \item $q_1^*$ максимизирует ожидаемый выигрыш фирмы 1 при равновесных выпусках фирмы 2 ($q_2^*(H)$ и $q_2^*(L)$):
    \[
    \max_{q_1} u_1 = \theta(1 - q_1 - q_2^*(H) - c_1)q_1 + (1 - \theta)(1 - q_1 - q_2^*(L) - c_1)q_1
    \]
    \item $q_2^*(H)$ максимизирует выигрыш фирмы 2 при высоких издержках $c_H$ и равновесном выпуске фирмы 1 ($q_1^*$):
    \[
    \max_{q_2} u_2 = (1 - q_1^* - q_2 - c_H)q_2
    \]
    \item $q_2^*(L)$ максимизирует выигрыш фирмы 2 при низких издержках $c_L$ и равновесном выпуске фирмы 1 ($q_1^*$):
    \[
    \max_{q_2} u_2 = (1 - q_1^* - q_2 - c_L)q_2
    \]
\end{itemize}

Условия первого порядка (FOC) для фирмы 2:
\begin{align*}
    q_2^*(H) &= \frac{1 - q_1^* - c_H}{2}, \\
    q_2^*(L) &= \frac{1 - q_1^* - c_L}{2}.
\end{align*}

Условие первого порядка для фирмы 1:
\[
q_1^* = \frac{1 - \theta q_2^*(H) - (1 - \theta) q_2^*(L) - c_1}{2}.
\]

Решение системы уравнений дает равновесные выпуски:
\begin{align*}
    q_2^*(H) &= \frac{1 - 2c_H + c_1}{3} + \frac{1 - \theta}{6}(c_H - c_L), \\
    q_2^*(L) &= \frac{1 - 2c_L + c_1}{3} - \frac{\theta}{6}(c_H - c_L), \\
    q_1^* &= \frac{1 - 2c_1 + \theta c_H + (1 - \theta) c_L}{3}.
\end{align*}
\end{example}

\begin{example}[Ценность информации в стратегическом взаимодействии]
В задачах индивидуальной оптимизации дополнительная информация не может ухудшить положение лица, принимающего решение (ее всегда можно проигнорировать). В стратегических играх это свойство может не выполняться.

Рассмотрим байесову игру двух лиц с двумя равновероятными состояниями $\omega_1$ и $\omega_2$. Пусть $0 < \varepsilon < 1/2$. Матрицы выигрышей:
\begin{table}
\centering
\begin{NiceTabular}{c|ccc}
\multicolumn{4}{c}{\textbf{Состояние $\omega_1$}} \\
\toprule
& \textit{L} & \textit{M} & \textit{R} \\
\midrule
$T$ & $1, 2\varepsilon$ & $1, 0$ & $1, 3\varepsilon$ \\
$B$ & $2, 2$ & $0, 0$ & $0, 3$ \\
\bottomrule
\end{NiceTabular}
\quad
\begin{NiceTabular}{c|ccc}
\multicolumn{4}{c}{\textbf{Состояние $\omega_2$}} \\
\toprule
& \textit{L} & \textit{M} & \textit{R} \\
\midrule
$T$ & $1, 2\varepsilon$ & $1, 3\varepsilon$ & $1, 0$ \\
$B$ & $2, 2$ & $0, 3$ & $0, 0$ \\
\bottomrule
\end{NiceTabular}
\end{table}

\textbf{Вариант 1: Никто не знает состояние мира.}
Игроки максимизируют ожидаемые выигрыши (игра со средними платежами):
\begin{table}
\centering
\begin{NiceTabular}{c|ccc}
\toprule
& \textit{L} & \textit{M} & \textit{R} \\
\midrule
$T$ & $1, 2\varepsilon$ & $1, 1.5\varepsilon$ & $1, 1.5\varepsilon$ \\
$B$ & $2, 2$ & $0, 1.5$ & $0, 1.5$ \\
\bottomrule
\end{NiceTabular}
\end{table}
По итеративному исключению доминируемых стратегий единственное равновесие Байеса-Нэша — профиль $(B, L)$. Ожидаемый платеж игрока 2 равен 2.

\textbf{Вариант 2: Игрок 2 знает состояние мира.}
Игрок 2 обладает точной информацией о матрице. В каждой матрице строго доминирующая стратегия для игрока 2:
\begin{align*}
    a_2^*(\omega_1) &= R, \\
    a_2^*(\omega_2) &= M.
\end{align*}
Игрок 1 сравнивает свои ожидаемые полезности с учетом оптимальной реакции информированного оппонента:
\begin{align*}
    \mathrm{E}u_1(T, a_2^*(\omega_1), a_2^*(\omega_2)) &= \frac{1}{2} \cdot 1 + \frac{1}{2} \cdot 1 = 1, \\
    \mathrm{E}u_1(B, a_2^*(\omega_1), a_2^*(\omega_2)) &= \frac{1}{2} \cdot 0 + \frac{1}{2} \cdot 0 = 0.
\end{align*}
Следовательно, оптимальная стратегия первого игрока: $a_1^* = T$. Ожидаемый платеж игрока 2 теперь равен $3\varepsilon$. Поскольку $\varepsilon < 1/2$, то $3\varepsilon < 2$. Получение информации ухудшило результат игрока 2 из-за изменения поведения неинформированного оппонента.
\end{example}

\begin{example}[Производство общественного блага при неполной информации]
Два игрока ($i = 1, 2$) одновременно принимают бинарное решение $\{0, 1\}$ о внесении вклада в создание общественного блага (1 — вкладываться, 0 — нет). Если хотя бы один игрок делает вклад, благо создается, и оба получают выгоду 1. Если никто не вкладывается, благо не создается, выигрыши равны 0. 

Издержки участия для игрока $i$ составляют $c_i$. Эти издержки известны только самому игроку (тип игрока). Общее знание состоит в том, что издержки распределены независимо согласно непрерывной и строго возрастающей функции распределения $P(\cdot)$ на носителе $[\underline{c}, \overline{c}]$, где $\underline{c} < 1 < \overline{c}$. 

\begin{table}
\centering
\begin{NiceTabular}{l|cc}
\toprule
& \textit{Contribute} & \textit{Don't} \\
\midrule
\textit{Contribute} & $1 - c_1, 1 - c_2$ & $1 - c_1, 1$ \\
\textit{Don't}      & $1, 1 - c_2$       & $0, 0$ \\
\bottomrule
\end{NiceTabular}
\end{table}

Чистая стратегия игрока $i$ — это функция $s_i(c_i) : [\underline{c}, \overline{c}] \to \{0, 1\}$. Выигрыш игрока $i$:
\[
u_i(s_1, s_2, c_i) = \max\{s_1, s_2\} - c_i s_i.
\]
Равновесие Байеса-Нэша задается парой функций $s_1^*(\cdot), s_2^*(\cdot)$, при которых стратегия $s_i^*(c_i)$ максимизирует ожидаемую полезность $\mathrm{E}_{c_j} u_i(s_i, s_j^*(c_j), c_i)$.

Обозначим через $z_j = P(s_j^*(c_j) = 1)$ "агрегированную" вероятность того, что игрок $j$ вносит вклад. Ожидаемые полезности действий игрока $i$:
\begin{align*}
    \mathrm{E}_{c_j} u_i(1, s_j^*(c_j), c_i) &= 1 - c_i \\
    \mathrm{E}_{c_j} u_i(0, s_j^*(c_j), c_i) &= z_j
\end{align*}

Оптимальный выбор игрока $i$: $s_i^*(c_i) = 1$, если $c_i < 1 - z_j$, и $s_i^*(c_i) = 0$, если $c_i > 1 - z_j$. Это означает, что равновесные стратегии являются \textit{монотонными} (пороговыми): игроки вкладываются при издержках ниже определенного значения $c_i^*$. Следовательно, интервал участия для игрока $i$ есть $[\underline{c}, c_i^*]$.

Поскольку $z_j = P(\underline{c} < c_j < c_j^*) = P(c_j^*)$, равновесный порог $c_i^*$ должен удовлетворять условию $c_i^* = 1 - P(c_j^*)$. Для симметричного равновесия ($c_1^* = c_2^* = c^*$) получаем уравнение:
\[
c^* = 1 - P(c^*)
\]

Рассмотрим пример с равномерным распределением $P(x) = x/2$ на отрезке $[0, 2]$. Уравнение принимает вид $c^* = 1 - c^*/2$, откуда $c^* = 2/3$.
Неэффективность равновесия: игрок не вкладывается, если его издержки лежат в интервале $(2/3, 1)$. Несмотря на то, что издержки меньше выгоды от создания блага ($c_i < 1$), игрок уклоняется от участия. С вероятностью $1 - P(2/3) = 2/3$ второй игрок также не вложится, и общественно полезное благо не будет произведено.

\begin{remark}
Если минимально возможные издержки достаточно высоки ($\underline{c} \ge 1 - P(1)$), в игре также существуют два асимметричных равновесия, в которых один игрок никогда не вкладывается, а другой всегда выступает единственным спонсором блага для любых $c \le 1$.
\end{remark}
\end{example}
